% =====================================================================
%  MCOT - Mise en Cohérence des Objectifs du TIPE (format SCEI / CPGE)
%  Sujet : optimisation du profil d'épaisseur d'un écran en verre trempé.
%
%  Respect des limites officielles de mots :
%    - Motivation .................... 50 mots max
%    - Ancrage au thème .............. 50 mots max
%    - Problématique ................. 50 mots max
%    - Objectifs ..................... 100 mots max
%    - Bibliographie commentée ....... 650 mots max
%  (les comptes approximatifs sont indiqués en commentaire % ~N mots)
%
%  Les MOTS TECHNIQUES sont mis en gras (\textbf{...}).
% =====================================================================
\documentclass[11pt,a4paper]{article}
\usepackage[utf8]{inputenc}
\usepackage[T1]{fontenc}
\usepackage[french]{babel}
\usepackage[margin=2.3cm]{geometry}
\usepackage{amsmath,amssymb}
\usepackage{enumitem}
\usepackage{titlesec}
\usepackage[hidelinks]{hyperref}

\titleformat{\section}{\large\bfseries\scshape}{}{0pt}{}
\setlength{\parindent}{0pt}
\setlength{\parskip}{4pt}

\begin{document}

% ---------------------------------------------------------------------
\begin{center}
{\Large\bfseries Mise en Cohérence des Objectifs du TIPE (MCOT)}\\[6pt]
{\large\bfseries Optimisation du profil d'épaisseur d'un écran protecteur\\
de smartphone en verre trempé}\\[8pt]
\small NOM Prénom \quad--\quad N\textsuperscript{o} candidat : \dots\dots \quad--\quad Filière MP/MPI \quad--\quad Année 20\dots--20\dots
\end{center}
\hrule
\bigskip

% ---------------------------------------------------------------------
\section{Positionnement thématique}
\textbf{Physique} (mécanique du solide, \textbf{résistance des matériaux}),
\textbf{Informatique pratique} (simulation numérique),
\textbf{Mathématiques} (probabilités).

\section{Mots-clés}
\begin{center}\small
\begin{tabular}{ll}
\textbf{Français} & \textbf{Anglais} \\[2pt]
verre trempé & tempered glass \\
flexion d'une poutre & beam bending \\
contrainte de rupture & fracture stress \\
simulation de Monte-Carlo & Monte-Carlo simulation \\
optimisation sous contrainte & constrained optimization \\
\end{tabular}
\end{center}

% --------------------------- 50 mots --------------------------------
\section{Motivation}
% ~50 mots
Chaque année, des millions d'écrans de smartphone se brisent lors d'une
chute. Passionné par la \textbf{mécanique} et la \textbf{physique des
matériaux}, j'ai voulu comprendre pourquoi la \textbf{rupture} s'amorce
aux \textbf{coins}, et tester une idée simple : une
\textbf{épaisseur variable} du \textbf{verre trempé} pourrait-elle réduire ce
risque tout en préservant la \textbf{tactilité} ?

% --------------------------- 50 mots --------------------------------
\section{Ancrage au thème de l'année}
% ~50 mots  (adapter le mot du thème officiel)
Le sujet s'inscrit dans le thème \emph{«~\dots(thème de l'année)\dots~»} : une
chute provoque une \textbf{transition} brutale de l'\textbf{énergie cinétique}
en \textbf{contrainte de flexion}, jusqu'à la \textbf{rupture fragile}.
Optimiser le \textbf{profil d'épaisseur} revient à chercher, par une
\textbf{boucle} de \textbf{simulation}, la configuration qui repousse ce
\textbf{seuil critique}.

% --------------------------- 50 mots --------------------------------
\section{Problématique retenue}
% ~50 mots
Quelle \textbf{répartition spatiale de l'épaisseur} d'un écran en
\textbf{verre trempé} minimise la \textbf{probabilité de rupture} lors d'une
chute, sous \textbf{contrainte de tactilité} (épaisseur centrale
$\le 0{,}5$~mm) et de \textbf{masse} ? On modélise l'écran par une
\textbf{poutre en flexion} et on estime cette probabilité par
\textbf{simulation de Monte-Carlo}.

% --------------------------- 100 mots ------------------------------
\section{Objectifs du TIPE}
% ~70 mots
\begin{enumerate}[leftmargin=1.6em,itemsep=1pt]
  \item Établir un \textbf{modèle mécanique} reliant l'épaisseur locale $e$ à la
        \textbf{contrainte maximale} de traction, et démontrer la loi
        $\sigma_{\max}\propto e^{-1/2}$.
  \item Valider cette dynamique par une \textbf{analogie électromécanique}
        (\textbf{circuit RLC}) mesurable à l'\textbf{oscilloscope}.
  \item Estimer par \textbf{Monte-Carlo} la \textbf{probabilité de rupture} de
        plusieurs \textbf{profils} (uniforme, parabolique, exponentiel).
  \item Déterminer le \textbf{profil optimal} sous contraintes de tactilité et
        de masse.
\end{enumerate}

% --------------------------- 650 mots ------------------------------
\section{Bibliographie commentée}
% ~640 mots
La compréhension du problème repose d'abord sur la \textbf{théorie des plaques
et des poutres}. L'ouvrage de référence de Timoshenko~[1] établit l'\textbf{équation
de la flexion} et introduit la \textbf{rigidité de flexion}
$D = E\,e^{3}/12(1-\nu^{2})$, dont la dépendance cubique en l'épaisseur est
l'ingrédient central de notre étude. Pour rester au niveau du programme, nous
réduisons la plaque à une \textbf{poutre sur deux appuis}, dont la
\textbf{raideur} $k = 48EI/L^{3}$ conserve la même loi $k \propto e^{3}$ ; la
\textbf{contrainte de flexion} $\sigma = My/I$ y est élémentaire. Cette
simplification, discutée dans~[1], surestime légèrement la contrainte, ce qui
rend notre conclusion \textbf{conservative}. Cette réduction
\textbf{unidimensionnelle} introduit toutefois une limite : un écran étant
\textbf{rectangulaire}, la bonne variable n'est pas la distance au centre mais
la \textbf{proximité du bord}, afin que la surépaisseur reste \textbf{uniforme}
le long du contour ; [1] fournit le cadre pour discuter cet écart entre modèle
\textbf{1D} et \textbf{plaque 2D}.

Le choc lui-même relève de la \textbf{mécanique du contact}. Johnson~[2]
développe la \textbf{théorie de Hertz}, qui fournit la \textbf{durée de contact}
et la \textbf{force maximale} lors de l'impact d'un solide. La comparaison de
cette durée à la \textbf{période propre} de la plaque justifie
l'\textbf{approximation impulsionnelle} : le transfert de \textbf{quantité de
mouvement} est quasi instantané. Nous en retenons surtout une idée robuste --
une \textbf{fraction de l'énergie cinétique} de chute est convertie en
\textbf{énergie de flexion} -- qui évite le calcul détaillé du contact tout en
donne le bon \textbf{ordre de grandeur} de la fléche, donc de la contrainte. Le
\textbf{coefficient de restitution} et la répartition de l'énergie entre rebond,
onde sonore et déformation du téléphone expliquent pourquoi seule une faible
fraction sert réellement à fléchir le verre.

Le verre étant un matériau \textbf{fragile}, sa rupture est par nature
\textbf{statistique}. L'article fondateur de Weibull~[3] propose une \textbf{loi
de probabilité} de rupture dépendant de la contrainte et d'un \textbf{module de
Weibull}, qui rend compte de la dispersion liée aux \textbf{micro-défauts} de
surface. Cette approche éclaire le choix d'une \textbf{modélisation
probabiliste} : plutôt qu'un seuil déterministe unique, nous tirons le
\textbf{seuil de rupture} selon une loi de dispersion, ce qui justifie l'emploi
d'une \textbf{simulation de Monte-Carlo}. La trempe déplace ce \textbf{seuil
effectif} vers des valeurs élevées, mais les \textbf{bords}, moins bien traités,
demeurent les zones d'\textbf{amorçage} privilégiées, que la dispersion spatiale
de la loi permet de représenter.

Les propriétés spécifiques du \textbf{verre trempé} proviennent de~[4] et~[5].
Gy~[4] décrit la \textbf{trempe chimique par échange ionique}
($\mathrm{K}^{+}/\mathrm{Na}^{+}$), qui crée une \textbf{précontrainte de
compression de surface} : la rupture n'apparaît que lorsque la \textbf{traction
appliquée} dépasse cette précontrainte, ce qui élève fortement la résistance
effective. Les fiches techniques du \textbf{verre aluminosilicate}~[5]
fournissent les valeurs numériques utilisées (\textbf{module d'Young},
\textbf{coefficient de Poisson}, \textbf{masse volumique}, \textbf{contrainte de
compression de surface}), indispensables pour des résultats réalistes. Elles permettent aussi d'estimer la
\textbf{profondeur de couche comprimée} (DOL) et de discuter la
\textbf{fabricabilité} d'un \textbf{gradient continu} d'épaisseur, point critique
du sujet.

Enfin, la \textbf{méthode numérique} s'appuie sur les principes de la
\textbf{simulation de Monte-Carlo}~[6] : on tire un grand nombre de scénarios
aléatoires (hauteur de chute, \textbf{angle d'impact} selon une
\textbf{gaussienne tronquée}, point d'impact), on calcule pour chacun la
contrainte par le modèle de poutre, puis on compte la proportion d'événements
dépassant le seuil. La \textbf{loi des grands nombres} garantit la convergence
de l'estimateur de la \textbf{probabilité de rupture}. Cette méthode permet
aussi de comparer plusieurs \textbf{profils d'épaisseur} et de réaliser une
\textbf{optimisation sous contraintes} (tactilité, masse).

L'ensemble de ces sources dessine une démarche cohérente : [1] donne le modèle
mécanique, [2] justifie le traitement du choc, [3] fonde le critère de rupture
probabiliste, [4]--[5] ancrent le modèle dans des données matériau réelles, et
[6] fournit l'outil numérique. La convergence des disciplines --
\textbf{résistance des matériaux}, \textbf{physique du contact},
\textbf{probabilités} et \textbf{informatique} -- confère au sujet sa richesse
et sa faisabilité au niveau MP/MPI.

% ---------------------------------------------------------------------
\section{Références bibliographiques}
\small
\begin{enumerate}[label={[\arabic*]},leftmargin=2.2em,itemsep=1pt]
  \item S.~P.~Timoshenko, S.~Woinowsky-Krieger, \textit{Theory of Plates and
        Shells}, 2\textsuperscript{e}~éd., McGraw-Hill, 1959.
  \item K.~L.~Johnson, \textit{Contact Mechanics}, Cambridge University Press,
        1985.
  \item W.~Weibull, \textit{A Statistical Distribution Function of Wide
        Applicability}, Journal of Applied Mechanics, vol.~18, p.~293--297, 1951.
  \item R.~Gy, \textit{Ion exchange for glass strengthening}, Materials Science
        and Engineering~B, vol.~149, n\textsuperscript{o}~2, p.~159--165, 2008.
  \item Corning Incorporated, \textit{Gorilla Glass --- propriétés mécaniques du
        verre aluminosilicate} (fiches techniques constructeur).
  \item W.~L.~Dunn, J.~K.~Shultis, \textit{Exploring Monte Carlo Methods},
        Elsevier, 2011.
\end{enumerate}

\end{document}
